%! Setting Up the SVD — Unit 7, Lesson 7.0 — Student Packet Cover Sheet
\documentclass[10pt]{article}
\usepackage{linalg-article}
\usepackage{linalg-boxes}
\usepackage{ltablex}
\keepXColumns

\begin{document}

% Full-bleed burgundy banner
\begin{tikzpicture}[remember picture, overlay]
  \fill[burgundy] (current page.north west)
    rectangle ([yshift=-0.9in]current page.north east);
\end{tikzpicture}
\vspace{-0.8in}

\begin{center}
  {\color{white}\textbf{\LARGE Linear Algebra}}\\[4pt]
  {\color{blushmid}\large\textbf{Unit 7 \quad The Singular Value Decomposition}}\\[6pt]
  {\color{white}\textbf{\large Lesson 7.0 \quad Setting Up the SVD}}
\end{center}

\vspace{0.22in}
\namedateperiod
\vspace{0.18in}

\begin{learningtargetbox}
\small
\begin{enumerate}[leftmargin=1.5em, itemsep=5pt]
  \item \ldots{} explain that a matrix stretches the \textbf{unit circle} into an \textbf{ellipse}, whose half-widths are the \textbf{singular values} $\sigma_1\ge\sigma_2\ge0$.
  \item \ldots{} find a perpendicular input pair $\mathbf{v}_1,\mathbf{v}_2$ sent to a perpendicular output pair with $A\mathbf{v}_i=\sigma_i\mathbf{u}_i$, and read the singular values.
  \item \ldots{} state the \textbf{SVD} $A=U\Sigma V^{\top}$ and explain why \emph{every} matrix has one.
\end{enumerate}
\end{learningtargetbox}

\vspace{0.14in}

\begin{tocbox}
\small
\renewcommand{\arraystretch}{1.5}
\begin{tabularx}{\linewidth}{@{} c l X r @{}}
\rowcolor{blush}
\textbf{\#} & \textbf{Component} & \textbf{Description} & \textbf{Score} \\
\midrule
1 & Warm-Up        & Spiral review: matrix $\times$ vector, the perpendicular (dot-product) test, length \& unit vectors & \blank{1.2cm} \\
2 & Guided Notes   & The unit circle $\to$ ellipse; singular values $\sigma_i$; $A\mathbf{v}_i=\sigma_i\mathbf{u}_i$; $A=U\Sigma V^{\top}$ & \blank{1.2cm} \\
3 & Group Activity & Reading singular values off diagonal, swap, and symmetric matrices                      & \blank{1.2cm} \\
4 & Exit Ticket    & Quick check: read the ellipse, perpendicular in/out, and justify the idea               & \blank{1.2cm} \\
5 & Homework       & Reading $\sigma$'s, input vs.\ output frames, the $A^{\top}A$ preview, and a look ahead    & \blank{1.2cm} \\
\midrule
  & \multicolumn{2}{r}{\textbf{Total}} & \blank{1.2cm} \\
\end{tabularx}
\end{tocbox}

\vspace{0.10in}

\begin{remindbox}
\small A matrix stretches the \textbf{unit circle} into an \textbf{ellipse}. The half-widths of that
ellipse are the \textbf{singular values} $\sigma_1\ge\sigma_2\ge0$ --- the stretch factors --- and there
is always a perpendicular pair of input directions $\mathbf{v}_1,\mathbf{v}_2$ sent to a perpendicular
pair of output directions $\mathbf{u}_1,\mathbf{u}_2$, with $A\mathbf{v}_i=\sigma_i\mathbf{u}_i$.
Collecting these gives the \textbf{Singular Value Decomposition} $A=U\Sigma V^{\top}$: rotate, stretch,
rotate. Unlike eigenvectors, \emph{every} matrix --- of any shape --- has one, and the singular values
are always real and never negative. This is the tool behind image compression and data science.
\end{remindbox}

\end{document}
